\section{Sets, Elements and Axiom of Extension}

\subsection{What is a set?}
\indent Informally speaking; students in a classroom, pencils
in a case and books in a library are all simple examples
of sets.
However; in the set theory we do not define the word 
\textit{set}, because every theory need some elementary concepts.

Let's suppose that we had defined the concept of sets.
Then; we would need some other concepts, and for defining 
them, we would need some other concepts and so on...

Since this process cannot goes to infinity, we need some
elementary concepts that we accept them without any definition.
In set theory, there are two undefined concepts: One is \textit{set} 
which is an object and the other one is \textit{belonging} 
which is a relation between two sets.

If $x$ belongs to $A$, we say $x$ \textit{is an element of} $A$ and write 
\[ x \in A. \]
Similarly; if $x$ does not belong to $A$, we say $x$ \textit{is not an element of}
$A$ and write \[ x \not\in A. \]
In our notes, just for convenience, we will use small case letters for elements of 
sets which will be denoted by upper and more complex letters. For example:
\[x \in A\] and \[A \in \mathcal{F}.\]

\subsection{Subsets}
If for all $x$; if $x \in A$ then $x \in B$; then we say that $A$ is a \textit{subset} 
of $B$ and denote with \[A \subseteq B\] or \[B \supseteq A.\]
Similarly; if $A$ is not a subset of $B$, we say $A$ is not a subset of $B$
and write \[A \not\subseteq B.\]
Now, we can prove simple theorems including subsets.

\begin{theorem} Every set is subset of itself. \end{theorem}
\begin{proof}[Proof]
   Since every element in $A$ is an element of $A$, $A$ is an subset of $A$. 
\end{proof}

\begin{theorem} Suppose $A$ is a subset of $B$ and $B$ is a subset 
of $C$. Then, $A$ is a subset of $C$.\end{theorem}
\begin{proof}
    Suppose $x$ is an arbitrary element in A. Since $A \subseteq B$, $x \in B$.
    And since $B \subseteq C$, $x \in C$. Therefore, $A \subseteq C$
\end{proof}

\subsection{Set Equality}
We have defined the concept of subsets. We also need to determine when two sets
are \textit{equal}. Since equality is a meta-concept, we need an axiom to define equality
in sets.
\begin{axiom*}[Axiom of Extension]
    Two sets are equal if and only if they have the same elements.
\end{axiom*}
Or formally, \[\forall A \forall B [A = B \iff \forall x (x \in A \iff x \in B)].\]
Or simply we can say, \textit{sets are determined by their elements}.

We say that \textit{$A$ is a proper subset of $B$}, if $A$ is a subset of $B$,
but $A$ and $B$ are not equal; and denote with \[A \subset B.\]

\begin{theorem} Suppose $A$ is a subset of $B$ and $B$ is a subset of $A$.
Then, A and B are equal. \end{theorem}
\begin{proof}[Proof]
    Since $A \subseteq B$, every element in $A$ must be in $B$. And since
    $B \subseteq A$, every element in $B$ must be in $A$. Thus, they have
    same elements. Therefore, $A = B$.
\end{proof}

